% =============================================================================
% Bloque 9: Backup (slides B1–B9) — para Q&A
% =============================================================================

% --- B0a: Protocolo estadístico — justificación detallada ---
\begin{frame}{Backup B0a: ¿Por qué este protocolo estadístico?}
\scriptsize
Cada herramienta responde una pregunta \textbf{distinta} sobre el efecto:

\vspace{0.3em}
\begin{tabular}{@{}p{2.6cm}p{3.6cm}p{6.2cm}@{}}
\toprule
\textbf{Herramienta} & \textbf{Pregunta} & \textbf{Justificación} \\
\midrule
$t$-pareada       & ¿La diferencia es real?       & Configs emparejadas por \textit{dataset $\times$ clasificador $\times$ semilla}: elimina varianza entre configs. \\
\addlinespace[1pt]
Bonferroni        & ¿Y con muchos tests?          & 7 métodos vs.\ SMOTE $\Rightarrow$ $\alpha=0{.}05/7$. Conservador, blindado contra falsos positivos. \\
\addlinespace[1pt]
Bootstrap CI 95\% & ¿Qué rango es plausible?      & No asume normalidad. 10\,000 remuestreos: percentiles 2.5 y 97.5 estables al 3er decimal. \\
\addlinespace[1pt]
$d$ de Cohen      & ¿Qué tan grande es?           & Separa significancia \textbf{estadística} de \textbf{práctica}: con $n{=}525$, $p$ puede ser bajo aunque $\Delta$ sea trivial. \\
\addlinespace[1pt]
Win-rate          & ¿Qué tan consistente?         & Media $+2{.}25$pp con 50\% win = volátil. Con 82\% win = robusto. \\
\bottomrule
\end{tabular}

\vspace{0.3em}
\begin{insight}
\scriptsize Los \textbf{4 apuntan} en la misma dirección $\Rightarrow$ conclusión difícil de atacar.
\end{insight}
\end{frame}

% --- B0b: Interpretación de d de Cohen ---
\begin{frame}{Backup B0b: $d$ de Cohen — interpretación}
\small
\begin{columns}[T]

\begin{column}{0.45\textwidth}
\textbf{Fórmula (pareada):}
\[
d = \frac{\overline{\Delta\text{F1}}}{\text{SD}(\Delta\text{F1})}
\]

\vspace{0.4em}
\textbf{Convenciones (Cohen 1988):}\\[2pt]
\scriptsize
\begin{tabular}{@{}lll@{}}
\toprule
$d$ & Magnitud & Interpretación \\
\midrule
0.2 & pequeño  & apenas distinguible \\
0.5 & medio    & visible a ojo \\
0.8 & grande   & claramente diferente \\
1.0+ & muy grande & casi sin solapamiento \\
\bottomrule
\end{tabular}
\end{column}

\begin{column}{0.55\textwidth}
\setbeamercolor{block title}{fg=white,bg=uaigreen}
\begin{block}{Nuestros valores}
\scriptsize
\begin{tabular}{@{}lrl@{}}
\toprule
Configuración           & $d$    & Magnitud      \\
\midrule
Global (soft-weighted)  & 0.74   & medio-grande  \\
SVC (lineal)            & 1.00   & muy grande    \\
Ridge                   & 1.03   & muy grande    \\
MLP                     & 0.73   & medio-grande  \\
LogReg                  & 0.82   & grande        \\
\midrule
emotion 6-clases        & 1.24   & muy grande    \\
20news 10-shot          & 0.95   & muy grande    \\
\bottomrule
\end{tabular}
\end{block}

\vspace{0.3em}
\begin{insight}
\scriptsize Con $d{=}0.74$ y $p{<}0.0001$: la diferencia es \textbf{real} \emph{y} \textbf{práctica}.
\end{insight}
\end{column}

\end{columns}
\end{frame}


% --- B0c: Filtros geométricos — diseño completo ---
\begin{frame}{Backup B0c: Filtros geométricos --- arquitectura y diseño}
\scriptsize
\begin{columns}[T]

\begin{column}{0.50\textwidth}
\textbf{1. LOF (Local Outlier Factor)}
\begin{itemize}\setlength\itemsep{0pt}
  \item Score: densidad local relativa a $k$ vecinos.
  \item \textbf{Decisión binaria}: score $> \theta_{\text{LOF}}$.
  \item \texttt{novelty=True}: entrenado en reales, evaluado en sintéticos.
\end{itemize}

\vspace{0.4em}
\textbf{2. Cascade (4 niveles)}
\begin{itemize}\setlength\itemsep{0pt}
  \item L1: distancia euclidiana al centroide.
  \item L2: + similitud coseno al centroide.
  \item L3: + pureza KNN ($k{=}10$).
  \item L4: + confianza (dist.\ a centroide).
  \item Score combinado: \textbf{media geométrica}.
  \item Selección: \textbf{top-N por ranking} (sin umbral).
\end{itemize}
\end{column}

\begin{column}{0.50\textwidth}
\textbf{3. Combined (LOF $\land$ similarity)}
\begin{itemize}\setlength\itemsep{0pt}
  \item Debe pasar \textbf{ambos}: LOF $> 0$ \emph{y} sim $> 0.5$.
  \item Operador \textbf{AND} $\Rightarrow$ muy restrictivo.
\end{itemize}

\vspace{0.4em}
\textbf{4. EmbeddingGuided (coverage)}
\begin{itemize}\setlength\itemsep{0pt}
  \item Maximiza cobertura: distancia a vecinos.
  \item Penaliza si está demasiado cerca.
  \item Coverage 0.6 + Quality 0.4 (centroide).
\end{itemize}

\vspace{0.4em}
\textbf{5. None (control)}
\begin{itemize}\setlength\itemsep{0pt}
  \item Sin filtrado (selección aleatoria).
\end{itemize}
\end{column}

\end{columns}

\vspace{0.4em}
\begin{insight}
\scriptsize Cascade L1 + ponderación suave \textbf{top-N} es el ganador: distancia simple + ranking.
\end{insight}
\end{frame}

% --- B0d: Cascada L0-L4 — ablación de niveles ---
\begin{frame}{Backup B0d: Ablación de niveles --- ¿más filtros, mejor?}
\small
\textbf{Cascada con selección top-N (sin ponderación suave, binary mask):}

\vspace{0.4em}
\begin{columns}[T]

\begin{column}{0.55\textwidth}
\centering
\scriptsize
\begin{tabular}{@{}lrrl@{}}
\toprule
Nivel & $\Delta$F1 vs.\ SMOTE & Win-rate & Filtros activos \\
\midrule
L0 (control)               & $-0.52$            & 50.0\%   & ninguno (random) \\
\rowcolor{uaigreen!12}
\textbf{L1}                & \alert{$+1.94$}    & 91.7\%   & distancia \\
L2                         & $+1.94$            & 91.7\%   & + similitud \\
L3                         & $+1.95$            & 91.7\%   & + KNN purity \\
L4 (full)                  & $+1.92$            & 91.7\%   & + confianza \\
\bottomrule
\end{tabular}

\vspace{0.4em}
\textbf{Con ponderación suave (\textit{top-N}, $n{=}168$):}\\[2pt]
\scriptsize
\begin{tabular}{@{}lrr@{}}
\toprule
Nivel & $\Delta$F1 vs.\ SMOTE & Win-rate \\
\midrule
\rowcolor{uaigreen!12}
\textbf{L1}        & \alert{$+3.66$}  & 84.5\% \\
L2                 & $+3.60$          & 83.9\% \\
L4 (full)          & $+3.60$          & 83.3\% \\
\bottomrule
\end{tabular}
\end{column}

\begin{column}{0.42\textwidth}
\setbeamercolor{block title}{fg=white,bg=uaiblue}
\begin{block}{Por qué L1 gana}
\scriptsize
\begin{itemize}\setlength\itemsep{1pt}
  \item Las 4 métricas están \textbf{altamente correlacionadas} (distancia + coseno + KNN + centroide $\to$ todas miden ``cercanía al cluster'').
  \item La media geométrica de señales redundantes \textbf{no aporta} información.
  \item L1 es \textbf{robusto en 10-shot}: KNN/LOF degeneran con $<20$ vecinos.
\end{itemize}
\end{block}

\vspace{0.3em}
\begin{insight}
\scriptsize \textbf{Simple gana a complejo}: añadir filtros NO mejora, levemente perjudica.
\end{insight}
\end{column}

\end{columns}
\end{frame}

% --- B0e: Cómo se eligieron los umbrales ---
\begin{frame}{Backup B0e: Elección de hiperparámetros}
\scriptsize
\begin{columns}[T]

\begin{column}{0.55\textwidth}
\textbf{LOF: \texttt{n\_neighbors} y \texttt{threshold}}\\[2pt]
\begin{tabular}{@{}lrrr@{}}
\toprule
Param & $k=5$ & $k=10$ & $k=20$ \\
\midrule
$\Delta$F1 vs.\ SMOTE & \alert{$+2.43$} & $+0.16$ & $+0.05$ \\
Win-rate              & 90.0\%         & 70.0\% & 67.5\% \\
\bottomrule
\end{tabular}

\vspace{0.2em}
{\itshape Default sklearn $k{=}20$ falla en 10-shot (no hay 20 vecinos por clase). $k{=}5$ es el sweet spot.}

\vspace{0.5em}
\textbf{LOF: \texttt{threshold}}\\[2pt]
\begin{tabular}{@{}lrr@{}}
\toprule
Threshold & $\Delta$F1 & Win-rate \\
\midrule
$-0.50$       & $+1.29$ & 80.0\% \\
$-0.30$       & $+0.06$ & 68.3\% \\
$\phantom{-}0.00$ (default sklearn) & $+0.88$ & 75.6\% \\
\bottomrule
\end{tabular}
{\itshape $\theta{=}0$ es la frontera matemática inlier/outlier (\texttt{decision\_function}).}
\end{column}

\begin{column}{0.42\textwidth}
\textbf{Cascade: \texttt{k\_neighbors} para KNN purity}
\begin{itemize}\setlength\itemsep{0pt}
  \item Fijado $k{=}10$ (compromiso 10/25/50-shot).
  \item L1 (sin KNN) no es sensible a este parámetro.
\end{itemize}

\vspace{0.4em}
\textbf{Combined: \texttt{sim\_threshold}}\\[2pt]
\begin{tabular}{@{}lrr@{}}
\toprule
$\theta_{\text{sim}}$ & $\Delta$F1 & Win-rate \\
\midrule
$0.3$ & $-0.68$ & 55.0\% \\
$0.5$ & $-5.71$ & 41.7\% \\
$0.7$ & $-1.10$ & 71.7\% \\
\bottomrule
\end{tabular}
{\itshape Cualquier $\theta_{\text{sim}}{>}0$ combinado con AND $\Rightarrow$ catastrófico.}

\vspace{0.4em}
\textbf{Top-N: peso mínimo y temperatura}
\begin{itemize}\setlength\itemsep{0pt}
  \item Normalización: \texttt{minmax} (vs.\ rank, z-score).
  \item Temperatura $\tau{=}0.5$ (peaked, $w \propto s^{1/\tau}$).
  \item Peso mínimo $w_{\min}$: \textbf{poco impacto} ($\pm 0.05$pp).
\end{itemize}
\end{column}

\end{columns}

\vspace{0.3em}
\begin{insight}
\scriptsize Hiperparámetros del cascade son \textbf{poco sensibles}; los de LOF y combined son \textbf{críticos}.
\end{insight}
\end{frame}

% --- B0f: Tabla completa de filtros con costo ---
\begin{frame}{Backup B0f: Comparación completa de filtros geométricos}
\centering
\scriptsize
\setlength{\tabcolsep}{4pt}
\begin{tabular}{@{}lrrlrl@{}}
\toprule
Filtro                & $\Delta$F1 & Win & Selección & LLM & Notas \\
\midrule
\rowcolor{uaigreen!12}
\textbf{cascade L1 + top-N} & \alert{$+3{.}66$} & 84.5\% & ranking top-$N$ & 3$\times$   & \textbf{Ganador} \\
cascade L2 (top-N)         & $+3{.}60$ & 83.9\% & ranking top-$N$ & 3$\times$  & Similar a L1 \\
cascade L4 (top-N)         & $+3{.}60$ & 83.3\% & ranking top-$N$ & 3$\times$  & Redundante \\
cascade L1 (binary)        & $+1{.}94$ & 91.7\% & ranking top-$N$ & 3$\times$  & Sin pesos suaves \\
cascade L3 (binary)        & $+1{.}95$ & 91.7\% & ranking top-$N$ & 3$\times$  & + KNN, no ayuda \\
\addlinespace[2pt]
LOF ($k{=}5$, $\theta{=}{-}0{.}5$) & $+2{.}43$ & 90.0\% & umbral (relajado)   & 6.2$\times$  & Segundo mejor \\
LOF ($k{=}20$, $\theta{=}0$)       & $+0{.}05$ & 67.5\% & umbral (estricto)   & 8$\times$    & Default sklearn falla \\
\addlinespace[2pt]
none (control)            & $+0{.}72$ & 62.9\% & sin filtrar          & 3$\times$    & LLM directo \\
\addlinespace[2pt]
\rowcolor{uaiorange!10}
embedding\_guided         & $-1{.}29$ & 55.0\% & cobertura greedy     & --           & Coverage-only falla \\
\rowcolor{uaiorange!10}
combined (LOF $\land$ sim) & $-3{.}28$ & 52.5\% & AND, $\sim\!23\%$ pasa & 32$\times$ & Demasiado restrictivo \\
\bottomrule
\end{tabular}

\vspace{0.2em}
{\scriptsize\itshape \textbf{Selección}: \emph{ranking top-$N$} toma los $N$ mejores del pool 3$\times$ generado (filtrado por ranking, no por umbral). \emph{Umbral}: rechaza candidatos bajo el score; el \% indica cuántos pasan. ``LLM'' = candidatos por muestra requerida.}

\vspace{0.2em}
\begin{insight}
\scriptsize \textbf{Tres principios confirmados:}
(1) \textit{Ranking} top-$N$ $>$ umbral duro.\hspace{0.5em}
(2) Distancia simple $\geq$ múltiples criterios.\hspace{0.5em}
(3) Operador \textbf{AND} destruye la señal.
\end{insight}
\end{frame}

% --- B0g: Hiperparámetros del LLM y ponderación suave ---
\begin{frame}{Backup B0g: Hiperparámetros LLM y ponderación suave}
\scriptsize
\begin{columns}[T]

\begin{column}{0.50\textwidth}
\setbeamercolor{block title}{fg=white,bg=uaiblue}
\begin{block}{LLM (Gemini 3 Flash)}
\scriptsize
\begin{tabular}{@{}lrl@{}}
\toprule
Param         & Valor & Justificación \\
\midrule
\texttt{temperature}      & $0.8$  & Diversidad / coherencia \\
\texttt{top\_p}, \texttt{top\_k} & default & SDK Gemini \\
\texttt{max\_tokens}      & $4000$ & $\sim 50$ muestras/llamada \\
back-trans.\ $T$          & $0.3$  & EN $\leftrightarrow$ ES conservador \\
\bottomrule
\end{tabular}

\vspace{0.3em}
\textbf{Ablation \texttt{temperature} (cascade L1):}\\[2pt]
\begin{tabular}{@{}lrr@{}}
\toprule
$T$ & sms\_spam & 20news \\
\midrule
$0.5$ & $-0.07$ & $+5.04$ \\
$0.7$ & $+1.44$ & $+4.16$ \\
$0.8$ & $-1.06$ & $+6.04$ \\
$0.9$ & $-0.18$ & $+6.25$ \\
$1.0$ & $+4.95$ & $+7.44$ \\
\bottomrule
\end{tabular}\\
{\itshape No monótono. $T{=}0.8$ por convención.}
\end{block}
\end{column}

\begin{column}{0.50\textwidth}
\setbeamercolor{block title}{fg=white,bg=uaibluedark}
\begin{block}{Ponderación suave (top-N)}
\scriptsize
\begin{tabular}{@{}lrl@{}}
\toprule
Param            & Valor    & Comentario \\
\midrule
$\tau$ ponder.   & $0.5$    & $w \propto s^{1/\tau}$, peaked \\
normalización    & minmax   & frente a rank, sigmoid \\
$w_{\min}$       & $0.0$    & sensibilidad $<\!0.2$pp \\
candidatos       & $3\times$ target & margen top-N \\
\bottomrule
\end{tabular}

\vspace{0.1em}
\textbf{Sweep Phase 2 (468 configs):}
\begin{tabular}{@{}lrr@{}}
\toprule
Param          & $\Delta$F1 & Win \\
\midrule
$\tau=0.5$     & \alert{$+4.10$} & 88.3\% \\
$\tau=1.0$     & $+3.67$         & 84.0\% \\
$\tau=2.0$     & $+3.41$         & 83.3\% \\
norm.\ minmax  & \alert{$+3.98$} & 89.6\% \\
norm.\ rank    & $+3.56$         & 83.3\% \\
\bottomrule
\end{tabular}
\end{block}
\end{column}

\end{columns}

\begin{insight}
\scriptsize $\tau{=}0.5$ \textbf{peaked} amplifica diferencias entre muestras buenas y mediocres $\Rightarrow$ máxima ganancia.
\end{insight}
\end{frame}

% --- B0h: Embeddings, clasificadores y validación ---
\begin{frame}{Backup B0h: Embeddings, clasificadores y protocolo}
\scriptsize
\begin{columns}[T]

\begin{column}{0.50\textwidth}
\setbeamercolor{block title}{fg=white,bg=uaiblue}
\begin{block}{Clasificadores (sklearn defaults)}
\scriptsize
\begin{tabular}{@{}ll@{}}
\toprule
Modelo & Hiperparámetros \\
\midrule
LogReg        & \texttt{max\_iter=1000} \\
SVC           & \texttt{kernel="linear"} \\
RandomForest  & \texttt{n\_estimators=100} \\
Ridge         & $\alpha = 1.0$ \\
MLP (custom)  & 128 hidden, 300 ep., lr $10^{-3}$ \\
\bottomrule
\end{tabular}

\vspace{0.2em}
{\itshape Sin tuning por dataset $\Rightarrow$ las diferencias son del \textbf{método de aumentación}.}
\end{block}

\vspace{0.3em}
\setbeamercolor{block title}{fg=white,bg=uaiblue}
\begin{block}{Protocolo}
\scriptsize
Seeds $\{42, 123, 456, 789, 1011\}$ \textperiodcentered\ $N$-shot $\{10, 25, 50\}$ \textperiodcentered\ LLM \% $\{5, 25, 50, 100\}$ \textperiodcentered\ macro F1 \textperiodcentered\ 50 sintét./clase \textperiodcentered\ cache MD5.
\end{block}
\end{column}

\begin{column}{0.50\textwidth}
\setbeamercolor{block title}{fg=white,bg=uaibluedark}
\begin{block}{Embedding model (504 runs)}
\scriptsize
\begin{tabular}{@{}lrrr@{}}
\toprule
Modelo & Dim & $\Delta$F1 & Win \\
\midrule
\rowcolor{uaigreen!12}
\textbf{mpnet-base} & 768  & \alert{$+2.74$} & 92.1\% \\
bge-large           & 1024 & $+2.04$         & 69.8\% \\
bge-small           & 384  & $+2.00$         & 76.2\% \\
e5-large            & 1024 & $+0.97$         & 68.3\% \\
\bottomrule
\end{tabular}

\vspace{0.2em}
{\itshape \textbf{Más grande $\neq$ mejor}: bge-small (384d) supera a e5-large (1024d).}
\end{block}
\end{column}

\end{columns}

\vspace{0.2em}
\begin{insight}
\scriptsize Hiperparámetros del \textbf{filtro} son robustos ($\pm 0.5$pp). Los del \textbf{embedding} y \textbf{clasificador} dominan la varianza.
\end{insight}
\end{frame}

% --- B1: Validación de SMOTE ---
\begin{frame}{Backup B1: SMOTE — validación de implementación}
\small
\begin{columns}[T]

\begin{column}{0.55\textwidth}
\textbf{Crítica anticipada del comité:}
\begin{quote}\itshape\scriptsize
``¿Su SMOTE \textit{dummy-binary} es equivalente al SMOTE estándar binary-mask?''
\end{quote}

\textbf{Validación:}
\begin{itemize}\setlength\itemsep{0pt}
  \item 27 comparaciones lado a lado.
  \item Diferencia máxima: $0{.}0000$ pp.
  \item F1 idénticos en todos los datasets.
\end{itemize}

\vspace{0.3em}
\begin{insight}
\scriptsize Línea base SMOTE \textbf{robusta y reproducible}.
\end{insight}
\end{column}

\begin{column}{0.45\textwidth}
\centering
\setbeamercolor{block title}{fg=white,bg=uaigreen}
\begin{block}{Resumen}
\scriptsize
\begin{tabular}{@{}lr@{}}
\toprule
Implementación   & F1 medio \\
\midrule
dummy-binary     & 0.7124 \\
binary-mask std. & 0.7124 \\
\midrule
Diferencia       & \textbf{0.0000} \\
\bottomrule
\end{tabular}
\end{block}
\end{column}

\end{columns}
\end{frame}

% --- B2: Seed asymmetry ---
\begin{frame}{Backup B2: Asimetría de semillas y caching}
\small
\textbf{Observación:} 80\% de combinaciones LLM $\times$ clasificador determinista
tienen \textbf{varianza cero} entre semillas (la generación está cacheada).

\vspace{0.4em}
\begin{columns}[T]

\begin{column}{0.55\textwidth}
\textbf{Validación rigurosa:}
\begin{itemize}\setlength\itemsep{0pt}
  \item Filtramos a clasificadores estocásticos\\(MLP, RandomForest).
  \item $\Delta$F1 = $+1{.}79$ pp.
  \item $p = 0{.}0015$, $d = 0{.}52$.
  \item Significancia se mantiene.
\end{itemize}
\end{column}

\begin{column}{0.45\textwidth}
\setbeamercolor{block title}{fg=white,bg=uaiamber}
\begin{block}{Implicancia}
\scriptsize
Los CIs subestiman la varianza real, pero las pruebas $t$ pareadas \textbf{siguen siendo válidas}.
\end{block}
\end{column}

\end{columns}
\end{frame}

% --- B3: Closed-loop (iterative prompting) ---
\begin{frame}{Backup B3: Closed-loop vs.\ pasada única}
\small
\textbf{Hipótesis exploratoria:} ¿mejora si re-prompteamos al LLM con feedback de los rechazos?

\vspace{0.5em}
\begin{tabular}{@{}lrr@{}}
\toprule
Estrategia                          & $\Delta$F1 vs.\ SMOTE & Win-rate \\
\midrule
Closed-loop (5 iteraciones)         & $+1{.}12$             & 52\% \\
Pasada única \alert{$\star$}        & $+2{.}08$             & 80\% \\
\bottomrule
\end{tabular}

\vspace{0.6em}
\textbf{¿Por qué falla closed-loop?}
\begin{itemize}\setlength\itemsep{0pt}
  \item Tasas de aceptación \textbf{decrecen} con iteraciones (-12.3 pp).
  \item Fallo dominante: \texttt{DISTANCE\_OUTLIER} (77\%).
  \item Cuello de botella: \textbf{el LLM}, no el prompt.
\end{itemize}

\vspace{0.4em}
\begin{insight}
\scriptsize Pasada única + filtro simple $>$ ingeniería iterativa de prompt.
\end{insight}
\end{frame}

% --- B4: Geometric vs Teacher ---
\begin{frame}{Backup B4: Filtrado geométrico vs.\ \textit{teacher}}
\small
\textbf{Comparación con un filtrado supervisado (teacher classifier):}

\vspace{0.5em}
\begin{tabular}{@{}lrrr@{}}
\toprule
Filtro & $\Delta$F1 vs.\ SMOTE & $p$ & Costo extra \\
\midrule
Geométrico (binary) & \alert{$+2{.}78$} & $<0{.}001$ & ninguno \\
Teacher (binary)    & $+2{.}43$         & $<0{.}001$ & entrenar clf adicional \\
\bottomrule
\end{tabular}

\vspace{0.6em}
\begin{insight}
\scriptsize El filtro geométrico es \textbf{más general} (no se especializa en separabilidad lineal) y \textbf{más eficiente} (no requiere entrenamiento extra).
\end{insight}
\end{frame}

% --- B5: Tabla completa de baselines modernos ---
\begin{frame}{Backup B5: Comparación con baselines modernos}
\small
\textbf{1.890 configuraciones, 10 métodos, 3 clasificadores:}

\vspace{0.5em}
\begin{tabular}{@{}lrrr@{}}
\toprule
Método                         & $\Delta$F1 vs.\ SMOTE & $p$ Bonf. & Win-rate \\
\midrule
\textbf{Soft-weighted (ours)}  & \alert{$+2{.}61$}     & $<0{.}0001$ & 88.9\% \\
\textbf{Binary-filter (ours)}  & $+2{.}39$             & $<0{.}0001$ & 79.4\% \\
No augmentation                & $+0{.}64$             & $<0{.}05$   & 56.0\% \\
Embedding mixup                & $-0{.}01$             & n.s.        & 51.0\% \\
Contextual BERT                & $-0{.}26$             & n.s.        & 48.0\% \\
T5 paraphrase                  & $-0{.}35$             & n.s.        & 47.0\% \\
Random oversample              & $-0{.}32$             & n.s.        & 47.0\% \\
EDA                            & $-0{.}49$             & $<0{.}05$   & 41.0\% \\
Back-translation               & $-0{.}50$             & $<0{.}05$   & 40.0\% \\
\bottomrule
\end{tabular}

\vspace{0.5em}
\begin{insight}
\scriptsize Todas las técnicas modernas \textbf{pierden} contra SMOTE; solo el filtrado geométrico gana.
\end{insight}
\end{frame}

% --- B6: Escasez de datos (movida desde Problema) ---
\begin{frame}{Backup B6: Escasez de datos --- un problema cotidiano}
\begin{columns}[T]

\begin{column}{0.5\textwidth}
\small
\textbf{Dominios cotidianos en clasificación de texto:}\\
Triaje clínico $\cdot$ tickets de soporte $\cdot$ spam $\cdot$ discurso de odio.

\vspace{0.6em}
\textbf{Por qué duele:}
\begin{itemize}\setlength\itemsep{3pt}
  \item Etiquetado humano = \alert{caro, lento y sesgado}.
  \item Las clases raras suelen ser \textbf{las que más importan}: urgencias, fraudes, diagnósticos.
  \item Pocos ejemplos $\Rightarrow$ el modelo \textbf{sobreajusta} y no generaliza a casos nuevos.
  \item Multiclase $+$ desbalance $\Rightarrow$ el rendimiento en clases raras colapsa.
\end{itemize}
\cita{Gopali 2024; Whitehouse 2023}
\end{column}

\begin{column}{0.5\textwidth}
\centering
\begin{tikzpicture}[every node/.style={font=\scriptsize}]
  \node[draw=uaiblue, rounded corners, fill=uaibluelight!15, minimum width=2.3cm, minimum height=1cm, text width=2.1cm, align=center] (a) at (-1.3,3) {\textbf{Hate speech}\\\tiny minoría: $<$30 ej.};
  \node[draw=uaiblue, rounded corners, fill=uaibluelight!15, minimum width=2.3cm, minimum height=1cm, text width=2.1cm, align=center] (b) at (1.3,3) {\textbf{Clínico}\\\tiny rare disease};
  \node[draw=uaiblue, rounded corners, fill=uaibluelight!15, minimum width=2.3cm, minimum height=1cm, text width=2.1cm, align=center] (c) at (-1.3,1.5) {\textbf{Soporte TI}\\\tiny producto nuevo};
  \node[draw=uaiblue, rounded corners, fill=uaibluelight!15, minimum width=2.3cm, minimum height=1cm, text width=2.1cm, align=center] (d) at (1.3,1.5) {\textbf{Spam}\\\tiny esquema emergente};
  \node[draw=uaiorange, very thick, rounded corners, fill=uaiorange!15, minimum width=4.8cm, minimum height=0.8cm, text width=4.6cm, align=center] at (0,0.2) {\textbf{Común:} 10--50 ejemplos por clase};
\end{tikzpicture}
\end{column}

\end{columns}
\end{frame}

% --- B7: Cuatro paradigmas (movida desde Estado del Arte) ---
\begin{frame}{Backup B7: Cuatro paradigmas, una década de evolución}
\centering
\begin{tikzpicture}[every node/.style={font=\scriptsize}]
  \draw[->, very thick, uaiblue!70] (0,0) -- (12,0) node[right, font=\scriptsize]{año};
  \foreach \x/\year in {0.5/2002, 3.5/2019, 6.5/2023, 9.5/2026}{
    \draw[uaiblue] (\x,-0.1) -- (\x,0.1);
    \node[below, font=\tiny] at (\x,-0.15) {\year};
  }

  \node[draw=uaiblue, fill=uaibluelight!15, rounded corners, text width=2.4cm, align=center,
        minimum height=1.2cm, font=\tiny] at (1.5, 1.5) (p1)
        {\textbf{Geométrico clásico}\\SMOTE, ADASYN, EDA};
  \draw[->, uaiblue!50] (0.5,0) -- (p1);

  \node[draw=uaiblue, fill=uaibluelight!15, rounded corners, text width=2.4cm, align=center,
        minimum height=1.2cm, font=\tiny] at (4.5, 1.5) (p2)
        {\textbf{Embeddings profundos}\\mixup, MixText, SMOTExT};
  \draw[->, uaiblue!50] (3.5,0) -- (p2);

  \node[draw=uaiblue, fill=uaibluelight!15, rounded corners, text width=2.4cm, align=center,
        minimum height=1.2cm, font=\tiny] at (7.5, 1.5) (p3)
        {\textbf{Generación LLM}\\GPT-3/4, Gemini, Claude};
  \draw[->, uaiblue!50] (6.5,0) -- (p3);

  \node[draw=uaiorange, very thick, fill=uaiorange!15, rounded corners, text width=2.4cm, align=center,
        minimum height=1.2cm, font=\tiny] at (10.5, 1.5) (p4)
        {\textbf{Híbridos + filtros}\\\textit{esta tesis}};
  \draw[->, uaiorange] (9.5,0) -- (p4);

  \node[font=\tiny, text width=2.6cm, align=center] at (1.5,-1) {Chawla 2002\\Wei \& Zou 2019};
  \node[font=\tiny, text width=2.6cm, align=center] at (4.5,-1) {Zhang 2018\\Chen 2020};
  \node[font=\tiny, text width=2.6cm, align=center] at (7.5,-1) {Suhaeni \& Yong 2023\\\textbf{+12.7\%} acc};
  \node[font=\tiny, text width=2.6cm, align=center] at (10.5,-1) {Cegin 2025, Springer 2025\\ImbLLM 2024, \textbf{esta tesis}};
\end{tikzpicture}

\vfill
\mensaje{Cada paradigma resuelve un problema y crea otro nuevo.}
\end{frame}

% --- B8: Una idea simple (movida desde Estado del Arte) ---
\begin{frame}{Backup B8: Una idea simple}
\vfill
\begin{center}
{\Large\itshape ¿Y si generamos muchos candidatos con un LLM\\[0.3em]
y dejamos que la \textbf{geometría} elija los útiles?}
\end{center}

\vspace{1.4em}
\begin{center}
\begin{tikzpicture}[every node/.style={font=\small, align=center}, node distance=1.0cm]
  \node[draw=uaiblue, fill=uaibluelight!20, rounded corners, very thick,
        text width=2.8cm, minimum height=1.4cm] (gen)
    {\textbf{Generar}\\\scriptsize candidatos\\con un LLM};

  \node[draw=uaiorange, fill=uaiorange!10, rounded corners, very thick,
        text width=2.8cm, minimum height=1.4cm, right=of gen] (filt)
    {\textbf{Filtrar}\\\scriptsize por geometría\\en \emph{embeddings}};

  \node[draw=uaigreen, fill=uaigreen!10, rounded corners, very thick,
        text width=2.8cm, minimum height=1.4cm, right=of filt] (use)
    {\textbf{Entrenar}\\\scriptsize con más datos\\útiles};

  \draw[->, very thick, uaigray] (gen) -- (filt);
  \draw[->, very thick, uaigray] (filt) -- (use);
\end{tikzpicture}
\end{center}

\vfill
\mensaje{Esta tesis evalúa cuándo esa idea funciona y cuándo no.}
\end{frame}

% --- B11: Las clases que más ganan (movida desde Resultados) ---
\begin{frame}{Backup B11: Las clases que más ganan --- las que partían perdiendo}
\begin{columns}[T]

\begin{column}{0.55\textwidth}
\centering
\scriptsize
\renewcommand{\arraystretch}{0.95}
\begin{tabular}{@{}llrrr@{}}
\toprule
\textbf{Dataset} & \textbf{Clase} & \textbf{SMOTE} & \textbf{Ours} & \textbf{$\Delta$} \\
\midrule
\rowcolor{uaiorange!12}\multicolumn{5}{l}{\textit{Clases con F1 base bajo}} \\
emotion       & \textbf{anger}         & 27.8 & 44.6 & \alert{\textbf{+16.9}} \\
20news        & \textbf{alt.atheism}   & 50.9 & 65.9 & \alert{\textbf{+15.0}} \\
emotion       & joy                    & 29.4 & 42.6 & \alert{+13.3} \\
hate\_speech  & \textbf{offensive}     & 73.3 & 86.3 & \alert{+13.0} \\
emotion       & fear                   & 48.8 & 59.0 & +10.2 \\
\midrule
\rowcolor{methodSMOTE!10}\multicolumn{5}{l}{\textit{Clases con F1 base alto}} \\
20news        & comp.graphics          & 90.8 & 91.3 & +0.5 \\
emotion       & sadness                & 41.3 & 42.3 & +1.0 \\
\bottomrule
\end{tabular}

\vspace{0.3em}
{\scriptsize\itshape 10-shot; 3 clasificadores lineales.}
\end{column}

\begin{column}{0.45\textwidth}
\small
\textbf{Patrón:} las clases \alert{difíciles inicialmente} son las que más ganan con filtrado geométrico.

\vspace{0.5em}
\begin{itemize}\setlength\itemsep{1pt}\scriptsize
  \item ``anger'' arranca en \textbf{27.8\%} (poco discriminable).
  \item Termina en \textbf{44.6\%}: $+16.9$ pp.
  \item ``comp.graphics'' parte ya en 90.8\%: poco margen, ganancia $+0.5$ pp.
\end{itemize}

\vspace{0.4em}
\begin{insight}
\scriptsize
Donde el clasificador \textbf{ya rinde bien}, el filtrado aporta poco.\\
Donde \textbf{tropieza}, el filtrado densifica la región correcta y \alert{refuerza la frontera}.
\end{insight}
\end{column}

\end{columns}

\mensaje{El método actúa justo donde el modelo más lo necesita.}
\end{frame}

% --- B10: Por qué estos datasets (movida desde Metodología) ---
\begin{frame}{Backup B10: ¿Por qué estos datasets de clasificación?}
\begin{columns}[T]

\begin{column}{0.50\textwidth}
\setbeamercolor{block title}{fg=white,bg=uaibluedark}
\begin{block}{\small Experimento principal --- 7 datasets}
\scriptsize
Buscan \textbf{cobertura en tres ejes}:
\begin{itemize}\setlength\itemsep{0pt}
  \item \textbf{Granularidad:} 2 $\to$ 20 clases.
  \item \textbf{Dominio:} spam, odio, noticias, emoción, ontología.
  \item \textbf{Estilo:} corto/largo, formal/informal, social/curado.
\end{itemize}
\vspace{0.4em}
\begin{itemize}\setlength\itemsep{0pt}
  \item \dataset{sms\_spam} (2): NLP canónico, frase corta.
  \item \dataset{hate\_speech\_davidson} (3): dominio sensible.
  \item \dataset{20newsgroups}, \dataset{ag\_news} (4): noticias.
  \item \dataset{emotion} (6): el ejemplo recurrente.
  \item \dataset{dbpedia14} (14): ontología Wikipedia.
  \item \dataset{20newsgroups\_20class} (20): máxima granularidad.
\end{itemize}
\end{block}
\end{column}

\begin{column}{0.50\textwidth}
\setbeamercolor{block title}{fg=white,bg=uaiorange}
\begin{block}{\small Escalabilidad --- 3 datasets adicionales}
\scriptsize
Responden la pregunta:\\[2pt]
\alert{¿el filtrado sigue funcionando con muchas clases?}
\vspace{0.4em}
\begin{itemize}\setlength\itemsep{0pt}
  \item \dataset{trec6} (6): preguntas TREC, mismo régimen de clases que el principal pero distinto dominio. Sirve de \textbf{control}.
  \item \dataset{banking77} (77): intenciones bancarias, granularidad fina con clases semánticamente cercanas.
  \item \dataset{clinc150} (150): intenciones de diálogo de asistentes virtuales. \textbf{Caso extremo}.
\end{itemize}
\vspace{0.4em}
Permiten medir \textbf{degradación del filtrado} al aumentar el número de clases.
\end{block}
\end{column}

\end{columns}

\mensaje{Diez datasets, dos preguntas: ¿qué tan bueno es el filtrado? ¿hasta dónde escala?}
\end{frame}

% --- B9: Dónde encaja esta propuesta (movida desde Estado del Arte) ---
\begin{frame}{Backup B9: Dónde encaja esta propuesta}
\begin{center}
\begin{tikzpicture}[scale=0.72, every node/.style={font=\scriptsize}, transform shape]
  \fill[uaiblue!25, opacity=0.7] (-1.4,0) circle (1.9);
  \fill[uaigreen!25, opacity=0.7] (1.4,0) circle (1.9);
  \fill[uaiamber!25, opacity=0.7] (0,-1.9) circle (1.9);

  \node[font=\bfseries\small, color=uaiblue, text width=2.4cm, align=center] at (-3.2,1.4) {Validación post-generación};
  \node[font=\bfseries\small, color=uaigreen!70!black, text width=2.4cm, align=center] at (3.2,1.4) {Ponderación adaptativa};
  \node[font=\bfseries\small, color=uaiamber!70!black, text width=2.6cm, align=center] at (0,-3.4) {Validación multi-tarea (NER)};

  \node[font=\tiny, color=uaiblue!70!black] at (-2.2,0) {LOF};
  \node[font=\tiny, color=uaigreen!60!black, text width=1.9cm, align=center] at (2.2,0) {Importance Sampling};
  \node[font=\tiny, color=uaiamber!70!black] at (0,-1.3) {MELM 2022};

  \node[draw=uaiorange, very thick, fill=white, rounded corners,
        font=\bfseries\scriptsize, text width=1.9cm, align=center,
        minimum height=0.8cm] at (0,-0.3) {Esta propuesta};
\end{tikzpicture}
\end{center}

\mensaje{Contribución: integrar las tres dimensiones por primera vez.}
\end{frame}
